\clearpage\corrigesChapitre{8}{Réactions chimiques et bilans}
\corrigesSouples
\begin{blocLivre}{cor-8-1}\begin{corrige}{1}{Lire une équation}
\begin{questions}
\item Réactifs : carbone et dioxygène ; produit : dioxyde de carbone.
\item Un atome de C et deux de O de chaque côté.
\end{questions}
\end{corrige}\end{blocLivre}
\begin{blocLivre}{cor-8-2}\begin{corrige}{2}{Ajuster l’eau}
\begin{questions}
\item $\ce{2 H2 + O2 -> 2 H2O}$.
\item Changer un indice changerait l’espèce chimique ; l’eau oxygénée n’est pas l’eau.
\end{questions}
\end{corrige}\end{blocLivre}
\begin{blocLivre}{cor-8-3}\begin{corrige}{3}{Un métal et un acide}
\begin{questions}
\item Le dihydrogène \ce{H2}.
\item À gauche $+2e$, à droite $+2e$ par jeu de coefficients ; la charge est conservée.
\end{questions}
\end{corrige}\end{blocLivre}
\begin{blocLivre}{cor-8-4}\begin{corrige}{4}{Un effet thermique}
\begin{questions}
\item Une transformation exothermique : de l’énergie est cédée au milieu réactionnel et au récipient.
\item La température, mesurée au cours du temps.
\end{questions}
\end{corrige}\end{blocLivre}
\begin{blocLivre}{cor-8-5}\begin{corrige}{5}{Un premier limitant}
\schema[corrige]{reaction}
\begin{questions}
\item $0,60/2=0,30$ ; $0,20/1=0,20$ : le dioxygène limite.
\item 0,40 mol d’eau se forment ; $0,60-0,40=\qty{0,20}{mol}$ de \ce{H2} restent.
\end{questions}
\end{corrige}\end{blocLivre}
\begin{blocLivre}{cor-8-6}\begin{corrige}{6}{Espèces spectatrices}
\begin{questions}
\item Les ions sodium et chlorure.
\item $\ce{H+ + HO- -> H2O}$.
\end{questions}
\end{corrige}\end{blocLivre}
\begin{blocLivre}{cor-8-7}\begin{corrige}{7}{Le méthane brûle}
\begin{questions}
\item $\ce{CH4 + 2 O2 -> CO2 + 2 H2O}$.
\item 0,300 mol de \ce{O2} sont consommées et 0,150 mol de \ce{CO2} produites.
\end{questions}
\end{corrige}\end{blocLivre}
\begin{blocLivre}{cor-8-8}\begin{corrige}{8}{Un solide reste}
\begin{questions}
\item Le carbonate de calcium est en excès puisqu’il n’a pas été entièrement consommé.
\item Les ions hydrogène de l’acide ont limité la réaction.
\end{questions}
\end{corrige}\end{blocLivre}
\begin{blocLivre}{cor-8-9}\begin{corrige}{9}{Détartrer avec des proportions}
\begin{questions}
\item $0,080/1=0,080$ contre $0,100/2=0,050$ : l’acide limite.
\item 0,050 mol de calcaire réagissent ; 0,030 mol restent ; 0,050 mol de \ce{CO2} se forment.
\end{questions}
\end{corrige}\end{blocLivre}
\begin{blocLivre}{cor-8-10}\begin{corrige}{10}{Conserver une masse totale}
\begin{questions}
\item $m=4,0+32,0=\qty{36,0}{g}$ par conservation de la masse.
\item De la vapeur ou un autre gaz peut quitter le système pesé ; il faut définir le système et contrôler les échanges.
\end{questions}
\end{corrige}\end{blocLivre}
\begin{blocLivre}{cor-8-11}\begin{corrige}{11}{Une proportion à ne pas confondre}
\begin{questions}
\item Les coefficients ne comparent pas les masses, car les entités n’ont pas la même masse.
\item Il faut deux moles de \ce{H2} pour une mole de \ce{O2}.
\end{questions}
\end{corrige}\end{blocLivre}
\begin{blocLivre}{cor-8-12}\begin{corrige}{12}{Comparer des essais thermiques}
\demarche{Comparer les quantités de réactif limitant dans des conditions identiques.}
\begin{questions}
\item Le nombre de transformations élémentaires du bilan double ; l’énergie libérée est alors proportionnelle à la quantité du limitant.
\item Les pertes thermiques, le temps de mesure et les échanges avec le récipient peuvent changer le résultat observé.
\end{questions}
\retiens{Une variation de température dépend de l’énergie libérée et des échanges avec l’environnement.}
\end{corrige}\end{blocLivre}

